\section{Diseño del sistema de subagentes}

\subsection{¿Por qué se necesita un subagente?}

\begin{importantbox}{Tres motivaciones principales}
\begin{enumerate}
    \item \textbf{Aislar el contexto}: evita que toda la información se acumule en la Context Window de un solo Master Agent y provoque un desbordamiento.
    \item \textbf{Ejecución concurrente}: varias tareas independientes pueden ser procesadas simultáneamente por distintos subagentes.
    \item \textbf{Entregables más extensos}: el Master Agent delega el trabajo pesado a los subagentes y solo se encarga del Reduce/consolidación final.
\end{enumerate}
\end{importantbox}

\subsection{Las cinco operaciones de subagente}

El sistema de subagentes de Codex está compuesto por cinco Function Tool:

\begin{enumerate}
    \item \textbf{Spawn Agents}: crea subagentes.
    \item \textbf{Send Inputs}: envía nuevas instrucciones a un subagente existente.
    \item \textbf{Wait}: bloquea la ejecución en espera de que uno o varios subagentes alcancen un estado final.
    \item \textbf{Close Agents}: cierra o da de baja explícitamente a un subagente.
    \item \textbf{Resume}: reanuda un subagente cerrado o finalizado, continuando el trabajo a partir de su contexto original.
\end{enumerate}

\begin{warningbox}{La comunicación con el subagente es unidireccional}
El subagente solo produce una salida \textbf{al terminar}. En el diseño estructural actual, el subagente no cuenta con un modo para interactuar de manera activa con el Master Agent: toda la comunicación va del Master Agent hacia el subagente.
\end{warningbox}

\subsection{Mecanismo de Agent Nickname}

El Codex Runtime mantiene un conjunto de Agent Nicknames, compuesto principalmente por nombres de científicos y académicos históricos. Al crear un subagente:
\begin{enumerate}
    \item Primero revisa si el archivo de configuración definió un nickname personalizado.
    \item En caso contrario, elige un nombre no ocupado del conjunto de nombres integrado.
    \item El nickname se usa para la presentación en el frontend; toda la interacción interna usa el Agent ID.
\end{enumerate}

\subsection{El valor de retorno de Spawn Agent y la percepción del contexto}

\begin{knowledgebox}{¿Cómo percibe el modelo el subagente que creó?}
La definición del Tool en el Function Calling no incluye una descripción del valor de retorno. Entonces, ¿cómo sabe el modelo qué subagente creó?

La respuesta está en que \textbf{el valor de retorno de Spawn Agent} (que incluye agent\_id y nickname) es reincorporado por el Runtime al contexto del modelo como Function Call Output. En su razonamiento posterior, el modelo puede leer naturalmente esta información del contexto y así establecer la asociación «qué agente creé y cuál es su ID».
\end{knowledgebox}

\subsection{La semántica de Wait Agent}

\begin{itemize}
    \item Wait no es un \texttt{sleep}, sino una \textbf{suscripción} al flujo de estados del agente; se parece más a una operación \texttt{join} / \texttt{await}.
    \item Bloquea el razonamiento posterior del Master Agent hasta que el subagente objetivo alcanza un estado final.
    \item Cuenta con un mecanismo de tiempo de espera (de hasta aproximadamente una hora).
    \item \textbf{Cuándo usarlo}: se activa solo cuando el siguiente paso del Master Agent depende necesariamente del resultado del subagente. Si todavía hay otras tareas que hacer, no conviene apresurarse a usar Wait.
\end{itemize}

\begin{importantbox}{Condiciones que activan Spawn Agent}
Codex establece \textbf{dos instrucciones estrictas} para el uso de Spawn Agent:
\begin{enumerate}
    \item Solo se crea un subagente cuando el usuario \textbf{solicita explícitamente} la asignación a subagentes o la ejecución concurrente.
    \item Excepción: si se trata de un \textbf{Research profundo e integral o un análisis sistemático} del Codebase, se puede crear un subagente sin necesitar el permiso del usuario.
\end{enumerate}
Por lo tanto, para probar el flujo de subagente, hay que mencionar explícitamente en la Query algo como «usar subagent para completarlo de forma concurrente».
\end{importantbox}

\subsection{Mecanismo de notificación al terminar el subagente}

Si el Master Agent no invoca Wait, el subagente no queda «desconectado» al terminar. En segundo plano, el Runtime inyecta automáticamente el mensaje de finalización del subagente \textbf{en calidad de User} en el contexto del Master Agent, con lo que se garantiza que la información no se pierda.

\subsection{Resumen de la sección}

El sistema de subagentes supera las limitaciones de una única Context Window mediante el aislamiento del contexto y la ejecución concurrente. Las cinco operaciones primitivas (Spawn / Send Inputs / Wait / Close / Resume) conforman la gestión completa del ciclo de vida del subagente.

